\chapter{The Major Continuous Distributions}

In the last chapter you learned the grammar of continuous random variables:
the density $f$, the cumulative distribution function $F$, the rule that
probabilities are \emph{areas} under the density, and the slogan that for a
continuous variable any single point has probability zero. Grammar without
vocabulary says nothing, though, so this chapter supplies the
vocabulary. We will meet the handful of continuous distributions that come up
again and again -- in physics, in queueing, in genetics, in finance, in
machine learning -- and we will do more than state their formulas. We will tell
the \emph{story} each one is built to model, watch the story force the density
into a particular shape, and then read the mean and variance off that shape.

There are four families to learn here, and they form a small, tightly knit
clan:

\begin{itemize}
  \item the \textbf{Uniform} -- ``I know it is somewhere in this interval and I
    have no further opinion'';
  \item the \textbf{Exponential} -- ``how long until the next event?'', the
    continuous cousin of the geometric, and the carrier of the famous
    \emph{memorylessness} property;
  \item the \textbf{Gamma} -- ``how long until the $k$-th event?'', the sum of
    independent exponentials, and the workhorse of waiting-time models;
  \item the \textbf{Beta} -- ``what is a plausible value for an unknown
    \emph{probability}?'', a distribution that lives on $[0,1]$ and that will
    become the secret hero of Bayesian inference.
\end{itemize}

We deliberately set the \emph{Normal} aside for its own chapter; it deserves
one. Everything here is connected: the exponential is a special gamma, the
gamma is a sum of exponentials, the beta is built from two gammas, and the
uniform is the simplest beta of all. These are not four formulas to memorize but
four chapters of one continuous story.

\section{The Uniform Distribution: Pure Ignorance}
\label{ch11:uniform}

\begin{intuitionbox}
Suppose I tell you a bus will arrive sometime in the next ten minutes, and that
is genuinely \emph{all} I know -- no schedule, no tendency to be early or late.
What is your belief about the arrival time? You should spread your belief evenly
over the interval $[0,10]$: no minute is more likely than any other. That even
spreading, made precise, is the uniform distribution. It is what ``I have no
reason to prefer any point over any other in this range'' looks like when
written as a density.
\end{intuitionbox}

\begin{definition}[Uniform distribution]
\label{ch11:def-uniform}
A random variable $X$ is \emph{uniform on $[a,b]$}, written
$X \sim \Unif(a,b)$, if it has the constant density
\[
  f(x) = \begin{cases} \dfrac{1}{b-a}, & a \le x \le b, \\[4pt] 0, & \text{otherwise.}\end{cases}
\]
\end{definition}

Why must the constant be $1/(b-a)$ and not something else? Because a density
must integrate to $1$. A flat rectangle of height $h$ sitting over an interval
of width $b-a$ has area $h(b-a)$, and setting that equal to $1$ forces
$h = 1/(b-a)$. The density does not care \emph{where} the interval sits, only
how \emph{wide} it is: a wider interval means a shorter rectangle, because the
same unit of total probability must be smeared over more room.

The CDF is just the running area of that rectangle. For $a \le x \le b$,
\[
  F(x) = \int_a^x \frac{1}{b-a}\,\dd t = \frac{x-a}{b-a},
\]
a straight line climbing from $0$ at $x=a$ to $1$ at $x=b$ (Figure~\ref{ch11:fig-uniform}).
The fraction $\tfrac{x-a}{b-a}$ is exactly ``how far along the interval are we?''
-- which is precisely the chance of landing to the left of $x$ when all points
are equally likely.

\begin{figure}[H]
\centering
\begin{tikzpicture}
\begin{axis}[
  width=11.5cm, height=5.2cm,
  xlabel={$x$}, ylabel={}, ymin=0, ymax=0.42,
  xmin=-0.5, xmax=5.5, samples=2,
  axis lines=left, legend pos=north east,
  legend style={font=\small}, title={Uniform$(1,4)$: density and CDF}]
% density: flat at 1/3 on [1,4]
\addplot[densityred, very thick, domain=1:4] {1/3};
\addplot[densityred, thick, dashed, domain=-0.5:1] {0};
\addplot[densityred, thick, dashed, domain=4:5.5] {0};
\addlegendentry{density $f(x)$}
% CDF: line from (1,0) to (4,1)
\addplot[distblue, very thick, domain=1:4] {(x-1)/3};
\addplot[distblue, thick, domain=-0.5:1] {0};
\addplot[distblue, thick, domain=4:5.5] {1};
\addlegendentry{CDF $F(x)$}
\end{axis}
\end{tikzpicture}
\caption{The uniform density is a flat rectangle; its CDF is the straight ramp
that accumulates that rectangle's area. Height $1/(4-1)=1/3$ makes the total
area equal $1$.}
\label{ch11:fig-uniform}
\end{figure}

\subsection{Mean and variance: the center and the spread of a rectangle}

By symmetry the average value of a uniform on $[a,b]$ must be the midpoint
$(a+b)/2$ -- the rectangle balances on its center. Let us confirm it and then
work out the variance.

\begin{align*}
  \E[X] &= \int_a^b x\cdot\frac{1}{b-a}\,\dd x
         = \frac{1}{b-a}\cdot\frac{b^2-a^2}{2}
         = \frac{(b-a)(b+a)}{2(b-a)} = \frac{a+b}{2}.
\end{align*}

For the variance, the cleanest route is $\Var(X) = \E[X^2] - (\E X)^2$. First,
\[
  \E[X^2] = \int_a^b x^2\cdot\frac{1}{b-a}\,\dd x
          = \frac{1}{b-a}\cdot\frac{b^3-a^3}{3}
          = \frac{a^2+ab+b^2}{3},
\]
using $b^3-a^3 = (b-a)(a^2+ab+b^2)$. Subtracting the square of the mean,
\[
  \Var(X) = \frac{a^2+ab+b^2}{3} - \frac{(a+b)^2}{4}
          = \frac{4(a^2+ab+b^2) - 3(a+b)^2}{12}
          = \frac{a^2 - 2ab + b^2}{12} = \frac{(b-a)^2}{12}.
\]

\begin{whybox}
Why $(b-a)^2/12$ and not $(b-a)^2/4$ or some other fraction? The variance can
depend only on the \emph{width} $w=b-a$, never on where the interval sits, since
shifting a distribution leaves its spread unchanged. Dimensionally it must be
$w^2$ times a pure number. That pure number is $1/12$. You can feel why it is
small and not, say, $1/4$: most of the probability is \emph{not} out at the
extreme ends, so the typical squared distance from the center is well under the
maximum possible $(w/2)^2 = w^2/4$. The average of those squared
distances works out to a third of that maximum, $w^2/12$. The standard
deviation is $w/\sqrt{12} \approx 0.289\,w$ -- about $29\%$ of the interval
width, a number worth remembering.
\end{whybox}

\begin{example}[Rounding error]
\label{ch11:ex-rounding}
When you round a measurement to the nearest unit, the error -- true value minus
rounded value -- is well modeled as $\Unif(-\tfrac12,\tfrac12)$: you have no
reason to think the discarded part favors any sub-part of the rounding window.
Its mean is $0$ (rounding is unbiased on average) and its variance is
$\frac{1^2}{12} = \frac{1}{12}$, so its standard deviation is
$1/\sqrt{12} \approx 0.289$ units. This single fact -- the
``\emph{quantization noise has variance $1/12$}'' rule -- underlies error
budgets in digital audio, sensors, and analog-to-digital converters.
\end{example}

\begin{aestheticbox}
The uniform is the distribution that admits the most ignorance: among
all distributions supported on $[a,b]$, it is the one with the largest entropy,
the one that assumes the least. We will make ``entropy'' precise in the
information-theory chapter, but the intuition is already here. To say
$X\sim\Unif(a,b)$ is to refuse to smuggle in any preference you do not actually
have.
\end{aestheticbox}

\begin{connectionbox}
The uniform is the gateway to \emph{simulating} every other distribution. If
$U\sim\Unif(0,1)$ and $F$ is any continuous, strictly increasing CDF, then
$X = F^{-1}(U)$ has CDF exactly $F$. This is the \emph{inverse-transform} or
\emph{probability integral transform} method, and it is why a single stream of
$\Unif(0,1)$ numbers from your computer can be reshaped into exponentials,
normals, anything. We prove it carefully in the chapter on transformations; for
now, hold the picture: the flat $U$ is the raw clay every other shape is
sculpted from. We will use it in a moment to build an exponential.
\end{connectionbox}

\section{The Exponential Distribution: How Long Until It Happens?}
\label{ch11:exponential}

\begin{intuitionbox}
A Geiger counter sits next to a lump of uranium. Clicks come at random. You
start a stopwatch. \emph{How long until the next click?} You cannot say exactly
-- it might be an instant, it might be a while -- but the longer you wait the
more it feels ``overdue.'' Or does it? The exponential distribution is the
mathematics of this waiting time, and it hides a shock: a decaying atom does not
age. The wait, no matter how long you have already waited, looks brand new. That
is memorylessness, and it is the most surprising thing in this chapter.
\end{intuitionbox}

\subsection{From the geometric to the exponential}

Recall the geometric distribution from the discrete chapter: flip a biased coin
with success probability $p$ on each trial; the number of trials until the first
success is geometric. The exponential is what the geometric \emph{becomes} when
trials happen continuously in time rather than one per tick.

Here is the limiting argument; doing it slowly shows where the formula comes
from. Chop time into tiny slices of width $\Delta t$. In
each slice, suppose an event happens with probability $\lambda\,\Delta t$, where
$\lambda > 0$ is a \emph{rate} (events per unit time), and slices are
independent. The chance of \emph{surviving} (no event) for a whole stretch of
length $t$, i.e.\ through $n = t/\Delta t$ slices, is
\[
  \P(T > t) = (1 - \lambda\,\Delta t)^{n}
            = \left(1 - \frac{\lambda t}{n}\right)^{n}
            \xrightarrow[\;n\to\infty\;]{} e^{-\lambda t},
\]
the famous limit $(1-x/n)^n \to e^{-x}$. So in the continuous limit the survival
function is $\P(T>t) = e^{-\lambda t}$. Everything else follows. The CDF is
$F(t) = 1 - e^{-\lambda t}$, and differentiating gives the density.

\begin{definition}[Exponential distribution]
\label{ch11:def-exponential}
A random variable $T \ge 0$ is \emph{exponential with rate $\lambda > 0$},
written $T \sim \Expo(\lambda)$, if it has density
\[
  f(t) = \lambda e^{-\lambda t}, \qquad t \ge 0,
\]
equivalently survival function $\P(T > t) = e^{-\lambda t}$ and
CDF $F(t) = 1 - e^{-\lambda t}$ for $t \ge 0$.
\end{definition}

The density starts at its highest point $\lambda$ at $t=0$ and decays
exponentially: short waits are the most likely, and very long waits are rare but
never impossible (Figure~\ref{ch11:fig-expo}). A bigger rate $\lambda$ means
events come faster, so the density is taller near zero and falls off more
steeply -- you do not wait long when things happen often.

\begin{figure}[H]
\centering
\begin{tikzpicture}
\begin{axis}[
  width=12cm, height=5.6cm,
  xlabel={$t$ (waiting time)}, ylabel={$f(t)$},
  ymin=0, ymax=2.1, xmin=0, xmax=4,
  samples=120, axis lines=left,
  legend style={font=\small}, legend pos=north east,
  title={Exponential densities $\lambda e^{-\lambda t}$ for three rates}]
\addplot[densityred, very thick, domain=0:4] {2*exp(-2*x)};
\addlegendentry{$\lambda=2$}
\addplot[distblue, very thick, domain=0:4] {1*exp(-1*x)};
\addlegendentry{$\lambda=1$}
\addplot[targetgreen, very thick, domain=0:4] {0.5*exp(-0.5*x)};
\addlegendentry{$\lambda=0.5$}
\end{axis}
\end{tikzpicture}
\caption{Every exponential density is a pure decay starting at height $\lambda$.
Faster rates (larger $\lambda$) crowd the probability toward short waits.}
\label{ch11:fig-expo}
\end{figure}

\subsection{Mean and variance: the rate sets the clock}

The mean wait should be the reciprocal of the rate: if events come at $3$ per
minute, you wait about $1/3$ of a minute on average. Here is the derivation. With
integration by parts (or by recognizing a gamma integral, which we do below),
\[
  \E[T] = \int_0^\infty t\,\lambda e^{-\lambda t}\,\dd t = \frac{1}{\lambda}.
\]
A slicker derivation uses the survival-function formula for nonnegative
variables, $\E[T] = \int_0^\infty \P(T>t)\,\dd t$:
\[
  \E[T] = \int_0^\infty e^{-\lambda t}\,\dd t = \frac{1}{\lambda}.
\]
For the second moment, $\E[T^2] = \int_0^\infty t^2 \lambda e^{-\lambda t}\,\dd t
= 2/\lambda^2$, so
\[
  \Var(T) = \frac{2}{\lambda^2} - \frac{1}{\lambda^2} = \frac{1}{\lambda^2},
  \qquad \SD(T) = \frac{1}{\lambda}.
\]

Notice something elegant: the standard deviation \emph{equals} the mean. The
exponential is exactly as spread out as it is large. This is a fingerprint --
whenever you see data whose standard deviation roughly equals its mean and which
is supported on the positive reals, suspect an exponential.

\begin{walkbox}
The mean of $\Expo(\lambda)$ is $1/\lambda$, but the \emph{median} -- the point
splitting probability in half -- is smaller. Solve $1 - e^{-\lambda t} = 1/2$ to
get $t = (\ln 2)/\lambda \approx 0.693/\lambda$. So the median is about $69\%$ of
the mean. Why is the median below the mean? Because the long right tail drags the
average up while leaving the halfway point behind. This right-skew is typical of
all positive waiting-time distributions.
\end{walkbox}

\subsection{Memorylessness: the property that defines the exponential}
\label{ch11:memoryless}

\begin{trapbox}
Here is the trap nearly everyone falls into. You have been waiting $10$ minutes
for the bus, and the wait is exponential. Surely the bus is now ``due'' -- your
remaining wait should be shorter than it was at the start? \textbf{No.} For an
exponential wait, the time you have already spent waiting tells you
\emph{nothing} about the time remaining. The distribution of your remaining wait
is identical to the distribution of a fresh wait. The clock does not remember.
\end{trapbox}

Let us prove it. The \emph{memoryless property} is the statement
\[
  \P(T > s + t \given T > s) = \P(T > t)
  \qquad \text{for all } s, t \ge 0.
\]
In words: given you have already survived past time $s$, the chance of surviving
an \emph{additional} $t$ is the same as the chance a brand-new wait survives $t$.

\begin{whybox}
The proof is one line, and it is the cleanest demonstration of why
exponentials are special. By the definition of conditional probability, and
because $\{T > s+t\} \subseteq \{T > s\}$,
\[
  \P(T > s+t \given T > s)
   = \frac{\P(T > s+t)}{\P(T > s)}
   = \frac{e^{-\lambda(s+t)}}{e^{-\lambda s}}
   = e^{-\lambda t} = \P(T > t).
\]
The whole thing turns on the algebraic miracle $e^{-\lambda(s+t)} =
e^{-\lambda s}\,e^{-\lambda t}$: the survival function \emph{factors}, so the
``already waited'' part cancels exactly. Memorylessness is nothing but the
multiplicative law of exponents wearing a probabilistic costume.
\end{whybox}

\begin{happeningbox}
What is really happening is that the exponential is the \emph{only} continuous
distribution with this property, and that is no coincidence -- it is forced. If
a positive variable satisfies $\P(T>s+t) = \P(T>s)\P(T>t)$ for all $s,t\ge0$,
then writing $g(t) = \P(T>t)$ we have $g(s+t)=g(s)g(t)$. The only measurable
solutions of this functional equation (Cauchy's exponential equation) that also
decrease from $1$ are $g(t) = e^{-\lambda t}$. So ``the wait does not age'' and
``the wait is exponential'' are two names for the same fact. Memorylessness is
not a quirky bonus feature of the exponential; it is its defining DNA.
\end{happeningbox}

\begin{paradoxbox}[The waiting-time paradox / the bus that is always late]
Memorylessness leads to a genuinely disorienting consequence. Buses arrive
according to a process whose gaps are exponential with mean $10$ minutes. You
show up at the stop at a random moment. How long do you expect to wait for the
next bus? Naively, ``I arrive in the middle of a gap, so about $5$ minutes.''
But memorylessness says: from the instant you arrive, the wait to the next bus
is a \emph{fresh} exponential with mean $10$ minutes -- you expect to wait
$10$, not $5$! And by the same logic the time \emph{since} the last bus also
averages $10$. So the gap you landed in averages $20$ minutes, twice the
average gap of $10$.

The resolution is \emph{length-biased sampling}: a randomly chosen \emph{moment}
is more likely to fall inside a long gap than a short one, simply because long
gaps cover more of the timeline. You are not sampling gaps uniformly; you are
sampling them in proportion to their length (Figure~\ref{ch11:fig-lengthbias}).
The gaps you tend to land in are the big ones, which is exactly why your wait feels worse than the ``average
gap'' suggests. This is not a flaw in the model; it is a real effect, and it is
why the line at the bus stop, or the queue you join, so often feels longer than
the advertised average.
\end{paradoxbox}

\begin{figure}[H]
\centering
\begin{tikzpicture}[scale=1.0]
% timeline
\draw[->, thick] (-0.3,0) -- (11.3,0) node[right]{time};
% bus arrival ticks at these positions (uneven gaps)
\foreach \x in {0,1.2,2.0,5.0,6.0,9.8,11.0} {
  \draw[thick] (\x,-0.12) -- (\x,0.12);
}
% label a couple of gaps
\draw[<->, gray] (2.0,-0.5) -- node[below,font=\scriptsize]{short} (5.0,-0.5);
\draw[<->, gray] (6.0,-0.5) -- node[below,font=\scriptsize]{long gap} (9.8,-0.5);
% random arrival instant, drawn falling into the long gap
\draw[densityred, very thick, ->] (8.0,1.0) -- (8.0,0.18);
\node[densityred, font=\scriptsize, above] at (8.0,1.0) {you arrive here};
\node[font=\scriptsize] at (5.4,0.45) {bus arrivals (ticks)};
\end{tikzpicture}
\caption{Length-biased sampling. Bus arrivals (vertical ticks) carve the timeline
into gaps of unequal length. A uniformly random arrival instant is more likely to
land inside a wide gap than a narrow one, simply because wide gaps occupy more of
the line -- so the gap you experience is longer than the average gap, and your
expected wait is larger than half the mean gap.}
\label{ch11:fig-lengthbias}
\end{figure}

\begin{historybox}
The exponential's memorylessness mirrors a physical fact discovered by Ernest
Rutherford and Hans Geiger around 1908: radioactive decay has no memory. An atom
that has not yet decayed is, statistically, indistinguishable from a freshly
minted one -- it has no internal ``odometer.'' This is why we describe
radioactivity by a \emph{half-life}: the time for half a sample to decay is the
median of an exponential, $(\ln 2)/\lambda$, and it is constant no matter how old
the sample is. The same constant-hazard idea reappears wherever failure is
``random'' rather than driven by wear -- which is why electronic components in
their useful-life phase, before wear-out sets in, are modeled exponentially.
\end{historybox}

\begin{connectionbox}[The hazard rate viewpoint]
Engineers think in terms of the \emph{hazard rate}: given survival to time $t$,
the instantaneous rate of failing right now, $h(t) = f(t)/(1-F(t))$. For the
exponential, $h(t) = \lambda e^{-\lambda t}/e^{-\lambda t} = \lambda$, a
\emph{constant}. Constant hazard \emph{is} memorylessness: the object is no more
likely to fail in the next instant after a year than on day one. Distributions
with \emph{increasing} hazard (the Weibull with shape $>1$) model wear-out --
things that age. Decreasing hazard models ``infant mortality.'' The exponential
sits exactly on the boundary, the unique no-aging case.
\end{connectionbox}

\begin{example}[Exponential by inverse transform]
\label{ch11:ex-inverse}
Now we cash in the connection from the uniform section. To simulate
$T\sim\Expo(\lambda)$, set $U\sim\Unif(0,1)$ and solve $F(T)=U$, i.e.\
$1 - e^{-\lambda T} = U$. Then $e^{-\lambda T} = 1-U$ and
\[
  T = -\frac{1}{\lambda}\ln(1-U).
\]
Since $1-U$ is also $\Unif(0,1)$, programmers often write
$T = -\tfrac{1}{\lambda}\ln U$. One uniform draw, one logarithm, and you have a
perfect exponential waiting time. This little formula runs inside every
discrete-event simulation -- network packet arrivals, customers at a teller,
particles in a detector.
\end{example}

\begin{example}[Minimum of exponentials: the race]
\label{ch11:ex-min}
Two independent components have lifetimes $T_1 \sim \Expo(\lambda_1)$ and
$T_2 \sim \Expo(\lambda_2)$. The system fails when the \emph{first} one does, so
its lifetime is $M = \min(T_1,T_2)$. When does $M$ exceed $t$? Only if
\emph{both} survive past $t$:
\[
  \P(M > t) = \P(T_1 > t)\,\P(T_2 > t) = e^{-\lambda_1 t}e^{-\lambda_2 t}
            = e^{-(\lambda_1+\lambda_2)t}.
\]
That is the survival function of an exponential with rate
$\lambda_1 + \lambda_2$. So the minimum of independent exponentials is again
exponential, and \emph{rates add}. Two independent Geiger counters pointed at the
same source behave like one counter with the combined click rate. This
``competing exponentials'' fact is the engine of the Poisson process and of
continuous-time Markov chains in later chapters.
\end{example}

\section{The Gamma Distribution: Waiting for the $k$-th Event}
\label{ch11:gamma}

\begin{intuitionbox}
The exponential answers ``how long until the \emph{next} click?'' Now ask: how
long until the \emph{tenth} click? You must wait for click one, then -- starting
fresh, by memorylessness -- for click two, and so on, ten times. The total wait
is a \emph{sum} of ten independent exponentials. That sum has its own shape: it
can no longer pile up at zero (ten waits can hardly all be tiny), so the density
rises from zero, humps in the middle, and trails off to the right. That humped,
right-skewed shape is the gamma distribution.
\end{intuitionbox}

\subsection{The gamma function, briefly}

The gamma distribution borrows its normalizing constant from the \emph{gamma
function}, the continuous extension of the factorial:
\[
  \Gamma(\alpha) = \int_0^\infty x^{\alpha-1} e^{-x}\,\dd x, \qquad \alpha > 0.
\]
Two facts are all we need. First, integration by parts gives the recursion
$\Gamma(\alpha+1) = \alpha\,\Gamma(\alpha)$, the same rule that defines the
factorial. Second, $\Gamma(1) = \int_0^\infty e^{-x}\dd x = 1$. Together these
force $\Gamma(n) = (n-1)!$ for positive integers $n$. So $\Gamma$ interpolates
the factorial to all positive reals -- it is the factorial with the gaps filled
in.

\begin{definition}[Gamma distribution]
\label{ch11:def-gamma}
A random variable $X \ge 0$ has the \emph{gamma distribution} with shape
$\alpha > 0$ and rate $\lambda > 0$, written $X\sim\Gam(\alpha,\lambda)$, if
\[
  f(x) = \frac{\lambda^{\alpha}}{\Gamma(\alpha)}\,x^{\alpha-1} e^{-\lambda x},
  \qquad x \ge 0.
\]
The factor $\lambda^{\alpha}/\Gamma(\alpha)$ is exactly what makes the area
equal $1$; you can check it by substituting $u = \lambda x$ in the integral and
recognizing the definition of $\Gamma(\alpha)$.
\end{definition}

Two pieces multiply to make the shape: the rising power $x^{\alpha-1}$ pulls the
density up away from zero, and the decaying $e^{-\lambda x}$ pulls it back down.
Their tug-of-war produces a hump. When $\alpha = 1$ the rising factor is
$x^0 = 1$ and the gamma collapses to a pure exponential -- the exponential is
just $\Gam(1,\lambda)$. As $\alpha$ grows the hump moves right and the shape
grows more symmetric (Figure~\ref{ch11:fig-gamma}); we will see in the
central-limit chapter that this is the gamma slowly turning into a bell.

\begin{figure}[H]
\centering
\begin{tikzpicture}
\begin{axis}[
  width=12cm, height=5.8cm,
  xlabel={$x$}, ylabel={$f(x)$},
  ymin=0, ymax=1.05, xmin=0, xmax=10,
  samples=140, axis lines=left,
  legend style={font=\small}, legend pos=north east,
  title={Gamma densities, rate $\lambda=1$, varying shape $\alpha$}]
% alpha=1 : exponential
\addplot[densityred, very thick, domain=0.001:10] {exp(-x)};
\addlegendentry{$\alpha=1$ (exponential)}
% alpha=2
\addplot[distblue, very thick, domain=0:10] {x*exp(-x)};
\addlegendentry{$\alpha=2$}
% alpha=3 : x^2 e^-x / 2
\addplot[targetgreen, very thick, domain=0:10] {x^2*exp(-x)/2};
\addlegendentry{$\alpha=3$}
% alpha=5 : x^4 e^-x / 24
\addplot[quotientorange, very thick, domain=0:10] {x^4*exp(-x)/24};
\addlegendentry{$\alpha=5$}
\end{axis}
\end{tikzpicture}
\caption{Increasing the shape $\alpha$ pushes the hump rightward and makes it
more symmetric. At $\alpha=1$ the gamma is exactly an exponential; for integer
$\alpha=k$ it is the wait for the $k$-th event.}
\label{ch11:fig-gamma}
\end{figure}

\begin{marginnotebox}
Beware: there are \emph{two} common parametrizations of the gamma. We use the
\emph{rate} $\lambda$ (so the exponential rate carries over cleanly). Others use
the \emph{scale} $\theta = 1/\lambda$, writing the density with $e^{-x/\theta}$
and dividing by $\theta^{\alpha}$. Software defaults differ -- always check
whether a library wants rate or scale, or you will be off by powers of
$\lambda$.
\end{marginnotebox}

\subsection{The gamma as a sum of exponentials}

The promised story -- ``wait for the $k$-th event'' -- is a theorem.

\begin{theorem}[Sum of i.i.d.\ exponentials is gamma]
\label{ch11:thm-sum-expo}
If $T_1,\dots,T_k$ are independent, each $\Expo(\lambda)$, then
$S_k = T_1 + \cdots + T_k \sim \Gam(k,\lambda)$.
\end{theorem}

\begin{whybox}
Here is the idea without the full convolution grind (we develop convolutions
properly in the sums chapter). The wait $S_k$ for the $k$-th event in a stream
arriving at rate $\lambda$ exceeds $t$ exactly when \emph{fewer than $k$} events
have arrived by time $t$. The number of events by time $t$ is Poisson with mean
$\lambda t$ (we build this carefully in the Poisson-process chapter). So
\[
  \P(S_k > t) = \P(\text{at most } k-1 \text{ events by } t)
              = \sum_{j=0}^{k-1} \frac{(\lambda t)^j}{j!}\,e^{-\lambda t}.
\]
Differentiate the CDF $1 - \P(S_k>t)$ with respect to $t$: the sum telescopes
beautifully (each term's derivative cancels against the next), leaving only
\[
  f_{S_k}(t) = \frac{\lambda^{k}}{(k-1)!}\,t^{k-1} e^{-\lambda t}
             = \frac{\lambda^{k}}{\Gamma(k)}\,t^{k-1} e^{-\lambda t},
\]
which is exactly $\Gam(k,\lambda)$. The shape parameter $\alpha = k$ \emph{is}
the number of exponentials you summed. The gamma is the exponential's story told
$k$ times over.
\end{whybox}

This instantly gives the mean and variance for free, by linearity and
independence:
\[
  \E[S_k] = k\cdot\frac{1}{\lambda} = \frac{k}{\lambda},
  \qquad
  \Var(S_k) = k\cdot\frac{1}{\lambda^2} = \frac{k}{\lambda^2}.
\]
And the formulas extend, by the same integral computation as before, to
\emph{any} real shape $\alpha$:
\[
  \boxed{\;\E[X] = \frac{\alpha}{\lambda}, \qquad
         \Var(X) = \frac{\alpha}{\lambda^2}\;}
  \qquad \text{for } X\sim\Gam(\alpha,\lambda).
\]
Both scale the obvious way: $\alpha$ exponentials stacked give $\alpha$ times the
mean and $\alpha$ times the variance, while a faster rate $\lambda$ shrinks both.

\begin{example}[How long until the 5th customer?]
\label{ch11:ex-fifth}
Customers walk into a shop at rate $\lambda = 2$ per hour, with exponential gaps.
How long until the fifth customer, and how variable is that wait? The wait is
$\Gam(5,2)$, so $\E = 5/2 = 2.5$ hours and $\Var = 5/4 = 1.25$, giving
$\SD = \sqrt{1.25}\approx 1.12$ hours. Compare a \emph{single} exponential wait,
whose SD ($0.5$ h) equals its mean ($0.5$ h): the five-event wait has SD that is
\emph{less than half} its mean. Summing independent waits averages out their
fluctuations -- the relative spread shrinks like $1/\sqrt{k}$. This shrinking is
the first whisper of the law of large numbers and the central limit theorem.
\end{example}

\begin{connectionbox}[The chi-squared distribution is a gamma]
You will meet the \emph{chi-squared} distribution constantly in statistics --
it governs sample variances and goodness-of-fit tests. It is secretly a gamma.
A chi-squared with $k$ degrees of freedom is exactly $\Gam(\alpha=k/2,
\lambda=1/2)$. In particular the square of a single standard normal is
$\Gam(1/2,1/2)$, and summing $k$ of them gives $\Gam(k/2,1/2)$ because shapes
add (next paragraph). When the statistics chapters invoke $\chi^2$, remember:
it is just a gamma in disguise, and everything you know here applies.
\end{connectionbox}

\begin{proposition}[Gammas with a common rate add]
\label{ch11:prop-gamma-add}
If $X\sim\Gam(\alpha,\lambda)$ and $Y\sim\Gam(\beta,\lambda)$ are independent
and share the same rate $\lambda$, then $X+Y\sim\Gam(\alpha+\beta,\lambda)$.
\end{proposition}

This is immediate from the sum-of-exponentials picture when $\alpha,\beta$ are
integers (concatenate the two stacks of waits), and holds for all positive
shapes by the convolution computation. \emph{Shapes add; the rate stays fixed.}
The exponential's ``rates add for minima'' (Example~\ref{ch11:ex-min}) and the
gamma's ``shapes add for sums'' are the two fundamental closure properties of
this family.

\section{The Beta Distribution: A Distribution Over Probabilities}
\label{ch11:beta}

\begin{intuitionbox}
Every distribution so far has lived on an unbounded stretch of the line. The
beta lives on the unit interval $[0,1]$ -- and that domain is the whole point.
What naturally lives in $[0,1]$? A \emph{probability}, or a \emph{proportion}, or
a \emph{fraction}. If you do not know the bias $p$ of a coin, or the click-through
rate of an ad, or the fraction of voters favoring a candidate, your uncertainty
about that unknown number is itself a distribution on $[0,1]$. The beta is the
natural language for ``what do I believe the probability is?'' -- a probability
distribution \emph{over} probabilities.
\end{intuitionbox}

\begin{definition}[Beta distribution]
\label{ch11:def-beta}
A random variable $X$ on $[0,1]$ has the \emph{beta distribution} with parameters
$\alpha > 0$ and $\beta > 0$, written $X\sim\Bet(\alpha,\beta)$, if
\[
  f(x) = \frac{1}{B(\alpha,\beta)}\,x^{\alpha-1}(1-x)^{\beta-1},
  \qquad 0 \le x \le 1,
\]
where the \emph{beta function} $B(\alpha,\beta) = \int_0^1
x^{\alpha-1}(1-x)^{\beta-1}\dd x$ is the constant that makes the area $1$.
\end{definition}

The shape is a tug-of-war between two factors. The term $x^{\alpha-1}$ pulls
mass toward $1$ (it grows as $x\to1$ when $\alpha>1$), while $(1-x)^{\beta-1}$
pulls mass toward $0$. Raising $\alpha$ favors large values; raising $\beta$
favors small values. When $\alpha=\beta$ the two pulls balance and the density
is symmetric about $1/2$. And when $\alpha=\beta=1$ both exponents vanish, the
density is the flat constant $1$, and the beta is exactly the uniform on
$[0,1]$: the uniform is the simplest beta, the one that says ``every probability
is equally plausible'' (Figure~\ref{ch11:fig-beta}).

\begin{figure}[H]
\centering
\begin{tikzpicture}
\begin{axis}[
  width=12cm, height=6cm,
  xlabel={$x$}, ylabel={$f(x)$},
  ymin=0, ymax=3.1, xmin=0, xmax=1,
  samples=140, axis lines=left,
  legend style={font=\small}, legend pos=north west,
  title={Beta densities for several $(\alpha,\beta)$}]
% Beta(1,1) = uniform
\addplot[gray, very thick, domain=0:1] {1};
\addlegendentry{$(1,1)$ uniform}
% Beta(2,2): 6 x (1-x)
\addplot[distblue, very thick, domain=0:1] {6*x*(1-x)};
\addlegendentry{$(2,2)$}
% Beta(5,2): x^4 (1-x) / B; B(5,2)=4!1!/6!=24/720=1/30 -> 30 x^4 (1-x)
\addplot[densityred, very thick, domain=0:1] {30*x^4*(1-x)};
\addlegendentry{$(5,2)$}
% Beta(2,5): mirror
\addplot[targetgreen, very thick, domain=0:1] {30*x*(1-x)^4};
\addlegendentry{$(2,5)$}
% Beta(0.5,0.5): U-shaped: 1/(pi sqrt(x(1-x)))
\addplot[quotientorange, very thick, domain=0.02:0.98] {1/(3.14159*sqrt(x*(1-x)))};
\addlegendentry{$(0.5,0.5)$}
\end{axis}
\end{tikzpicture}
\caption{The beta family covers an astonishing range of shapes on $[0,1]$:
flat $(1,1)$, hump-in-the-middle $(2,2)$, skewed toward $1$ $(5,2)$ or toward
$0$ $(2,5)$, and even U-shaped $(0.5,0.5)$ piling mass at both ends. This
flexibility is why it is the go-to model for an unknown proportion.}
\label{ch11:fig-beta}
\end{figure}

\subsection{The beta function and the beta--gamma bridge}

The beta function is not a new mystery; it is two gammas in a trench coat:
\[
  B(\alpha,\beta) = \frac{\Gamma(\alpha)\,\Gamma(\beta)}{\Gamma(\alpha+\beta)}.
\]

\begin{whybox}
This identity is where the four distributions of this chapter snap together.
Take independent $X\sim\Gam(\alpha,1)$ and $Y\sim\Gam(\beta,1)$. Consider the
\emph{fraction} $R = X/(X+Y)$ -- the share of the total contributed by $X$. It
is a number in $[0,1]$. A change of variables (done in full in the
transformations chapter) shows two clean facts: $R\sim\Bet(\alpha,\beta)$, and
$R$ is \emph{independent} of the total $X+Y\sim\Gam(\alpha+\beta,1)$. Matching the
normalizing constants in that computation is exactly the statement
$B(\alpha,\beta) = \Gamma(\alpha)\Gamma(\beta)/\Gamma(\alpha+\beta)$. So the beta
\emph{is} the distribution of ``what fraction of two competing gamma quantities
came from the first one.'' If $X$ is the time you wait for $\alpha$ red events
and $Y$ for $\beta$ blue events, the beta tells you what fraction of the total
waiting was red.
\end{whybox}

\subsection{Mean, variance, and how to read the parameters}

With the beta--gamma bridge, the mean is almost guessable: if $X$ contributes a
shape-$\alpha$ portion and $Y$ a shape-$\beta$ portion, the expected fraction
from $X$ should be $\alpha/(\alpha+\beta)$. It is:
\[
  \E[X] = \frac{\alpha}{\alpha+\beta},
  \qquad
  \Var(X) = \frac{\alpha\beta}{(\alpha+\beta)^2(\alpha+\beta+1)}.
\]
The mean is just the ratio of $\alpha$ to the total $\alpha+\beta$. The variance
deserves a closer look: write $s = \alpha+\beta$ for the ``total strength.'' Then
$\Var(X) = \E[X](1-\E[X])/(s+1)$. The numerator $\E[X](1-\E[X])$ is the spread
you would get from a single coin flip with that mean; the denominator $s+1$
\emph{shrinks} it. Larger $\alpha+\beta$ means a tighter, more confident
distribution clustered near its mean. This is the key to reading betas: the
\emph{ratio} $\alpha:\beta$ sets \emph{where} the belief sits, and the
\emph{sum} $\alpha+\beta$ sets \emph{how concentrated} it is.

\begin{example}[Reading a beta]
\label{ch11:ex-readbeta}
$\Bet(2,2)$ has mean $2/4 = 1/2$ and variance
$\frac{4}{16\cdot5} = 1/20 = 0.05$ ($\SD\approx0.224$): a gentle hump centered at
$1/2$, mild confidence. $\Bet(20,20)$ has the same mean $1/2$ but variance
$\frac{400}{1600\cdot41}\approx0.0061$ ($\SD\approx0.078$): same center, much
sharper -- ten times the strength, far more confident. And $\Bet(8,2)$ has mean
$8/10 = 0.8$: belief concentrated up near $80\%$. The two knobs do exactly what
their names promise.
\end{example}

\subsection{Conjugacy: a preview of Bayes}
\label{ch11:conjugacy}

Here is the property that makes the beta a superstar, and the reason it earns a
preview now even though Bayesian inference gets its own chapter later.

\begin{applicationbox}[The beta is the conjugate prior for a probability]
You want to learn the unknown success probability $p$ of a coin (or an ad, or a
drug). Before seeing data, encode your belief about $p$ as a beta:
$p \sim \Bet(\alpha,\beta)$. Now flip the coin $n$ times and observe $h$ heads
and $t = n-h$ tails. Bayes' rule (next part of the book) multiplies your prior
density by the likelihood of the data, which for $h$ heads and $t$ tails is
proportional to $p^{h}(1-p)^{t}$. So the updated (posterior) density is
proportional to
\[
  \underbrace{p^{\alpha-1}(1-p)^{\beta-1}}_{\text{prior}} \cdot
  \underbrace{p^{h}(1-p)^{t}}_{\text{likelihood}}
  = p^{(\alpha+h)-1}(1-p)^{(\beta+t)-1}.
\]
But that is just another beta! The posterior is
\[
  p \given \text{data} \;\sim\; \Bet(\alpha + h,\; \beta + t).
\]
The beta goes in, a beta comes out: this is \emph{conjugacy}. And the update rule
could not be more intuitive -- \textbf{add your heads to $\alpha$, add your tails
to $\beta$}. You can even read the prior as ``pretend experience'': $\Bet(\alpha,
\beta)$ behaves like having already seen $\alpha-1$ heads and $\beta-1$ tails
before the real data arrived. Starting from the uniform $\Bet(1,1)$ (``no prior
opinion'') and seeing $7$ heads in $10$ flips lands you at $\Bet(8,4)$, with
posterior mean $8/12 \approx 0.67$ -- pulled a touch toward $1/2$ from the raw
$0.7$, exactly the gentle regularization a sensible prior should provide.
Figure~\ref{ch11:fig-beta-update} shows a broad prior visibly sharpening into a
tighter posterior once data arrive.
\end{applicationbox}

\begin{figure}[H]
\centering
\begin{tikzpicture}
\begin{axis}[
  width=12cm, height=5.8cm,
  xlabel={$p$ (unknown click-through rate)}, ylabel={density},
  ymin=0, ymax=7, xmin=0, xmax=1,
  samples=140, axis lines=left,
  legend style={font=\small}, legend pos=north east,
  title={Bayesian update: a $\Bet(2,8)$ prior meets $12$ clicks in $40$ trials}]
% prior Beta(2,8): const = Gamma(10)/(Gamma(2)Gamma(8)) = 9!/(1!7!) = 362880/5040 = 72
\addplot[distblue, very thick, dashed, domain=0:1] {72*x*(1-x)^7};
\addlegendentry{prior $\Bet(2,8)$}
% posterior Beta(14,36): plot via exp of log-density for numerical stability.
% log normalizing constant = lnGamma(50)-lnGamma(14)-lnGamma(36) = 29.8774.
\addplot[densityred, very thick, domain=0.001:0.999, samples=160]
  {exp( 29.8774 + 13*ln(x) + 35*ln(1-x) )};
\addlegendentry{posterior $\Bet(14,36)$}
\end{axis}
\end{tikzpicture}
\caption{Data sharpens belief. The broad prior $\Bet(2,8)$ (mean $0.2$) is updated
by $12$ clicks and $28$ non-clicks into the much tighter posterior
$\Bet(14,36)$ (mean $0.28$). Conjugacy means the update is just ``add $12$ to
$\alpha$ and $28$ to $\beta$''; the posterior is pulled from the prior toward the
data proportion $12/40=0.30$.}
\label{ch11:fig-beta-update}
\end{figure}

\begin{happeningbox}
Why does conjugacy work? Look at the two ingredients. The beta density is
$p^{\,\square}(1-p)^{\,\square}$ in shape, and the coin-flip likelihood is also
$p^{\,\square}(1-p)^{\,\square}$ in shape. Multiplying functions of this form
just \emph{adds the exponents} -- the algebraic family $p^a(1-p)^b$ is closed
under multiplication. Conjugacy is not a deep coincidence; it is the statement
that the beta was \emph{designed} (or rather, discovered) to be the same shape as
the thing it multiplies. Counting heads and tails becomes literal arithmetic on
the parameters. This is why, in online systems that must update beliefs millions
of times a second -- A/B testing platforms, recommendation engines,
multi-armed-bandit ad servers -- the beta-binomial pair is the default: each new
click or non-click is one cheap increment of a counter.
\end{happeningbox}

\begin{connectionbox}
The same conjugacy idea, transplanted to waiting times, makes the gamma the
conjugate prior for an exponential or Poisson \emph{rate} $\lambda$: a gamma
prior on $\lambda$, updated by observed event counts and total time, yields a
gamma posterior, where you add the event count to the shape and the elapsed time
to the rate. So three of our four families -- beta for probabilities, gamma for
rates, and (in the next chapter) normal for means -- are precisely the
conjugate priors that make Bayesian updating into simple bookkeeping. The
continuous distributions you are learning here are the alphabet of Bayesian
statistics.
\end{connectionbox}

\section{The Family Portrait}
\label{ch11:portrait}

Step back and see the clan as a whole. Every relationship in this chapter is a
sentence about how one story becomes another.

\begin{itemize}
  \item \textbf{Uniform $=$ Beta$(1,1)$.} Flat ignorance is the simplest beta.
  \item \textbf{Exponential $=$ Gamma$(1,\lambda)$.} One waiting time is a
    one-step gamma.
  \item \textbf{Sum of $k$ Exponentials $=$ Gamma$(k,\lambda)$.} Stacking waits
    adds shapes.
  \item \textbf{Min of Exponentials $=$ Exponential, rates add.} Racing waits
    combine rates.
  \item \textbf{Gamma$+$Gamma (same rate) $=$ Gamma, shapes add.}
  \item \textbf{Beta $=$ Gamma$/(\text{Gamma}+\text{Gamma})$.} A beta is the
    fraction of a total contributed by one of two gammas.
  \item \textbf{Chi-squared $=$ Gamma$(k/2,1/2)$.} Statistics' favorite is a
    gamma in costume.
\end{itemize}

\begin{center}
\begin{tikzpicture}[
  node distance=12mm,
  box/.style={draw, rounded corners, align=center, font=\small,
              minimum width=24mm, minimum height=9mm},
  lab/.style={font=\scriptsize, fill=white, inner sep=1pt}]
\node[box, fill=blue!6]   (unif)  {Uniform$(0,1)$};
\node[box, fill=green!6, right=24mm of unif] (beta) {Beta$(\alpha,\beta)$};
\node[box, fill=red!6,   below=14mm of unif] (expo) {Exponential$(\lambda)$};
\node[box, fill=orange!8, right=24mm of expo] (gam)  {Gamma$(\alpha,\lambda)$};
\draw[-{Latex}] (unif) -- node[lab]{$\alpha{=}\beta{=}1$} (beta);
\draw[-{Latex}] (expo) -- node[lab]{$\alpha{=}1$} (gam);
\draw[-{Latex}] (gam.north) -- node[lab, right]{ratio $X/(X{+}Y)$} (beta.south);
\draw[-{Latex}] (expo.east) to[bend right=18]
   node[lab, below]{sum of $k$} (gam.west);
\end{tikzpicture}
\end{center}

\begin{aestheticbox}
Notice what organizes the whole picture: the exponent patterns. The gamma is
``power times decaying exponential,'' $x^{\alpha-1}e^{-\lambda x}$. The beta is
``power times power,'' $x^{\alpha-1}(1-x)^{\beta-1}$. Both are built from the
same atom, $x^{\text{power}}$, tamed on the right either by an exponential
(unbounded support) or by another power (bounded support). Seen this way the
families are not arbitrary; they are the two simplest ways to take a rising
power and make it integrable. That is the quiet aesthetic unity behind the zoo.
\end{aestheticbox}

\section{A Quick Reference}
\label{ch11:reference}

\begin{center}
\renewcommand{\arraystretch}{1.5}
\begin{tabular}{@{}lcccc@{}}
\toprule
Distribution & Support & Density & Mean & Variance \\
\midrule
$\Unif(a,b)$ & $[a,b]$ & $\dfrac{1}{b-a}$ & $\dfrac{a+b}{2}$ & $\dfrac{(b-a)^2}{12}$ \\
$\Expo(\lambda)$ & $[0,\infty)$ & $\lambda e^{-\lambda x}$ & $\dfrac{1}{\lambda}$ & $\dfrac{1}{\lambda^2}$ \\
$\Gam(\alpha,\lambda)$ & $[0,\infty)$ & $\dfrac{\lambda^\alpha}{\Gamma(\alpha)}x^{\alpha-1}e^{-\lambda x}$ & $\dfrac{\alpha}{\lambda}$ & $\dfrac{\alpha}{\lambda^2}$ \\
$\Bet(\alpha,\beta)$ & $[0,1]$ & $\dfrac{x^{\alpha-1}(1-x)^{\beta-1}}{B(\alpha,\beta)}$ & $\dfrac{\alpha}{\alpha+\beta}$ & $\dfrac{\alpha\beta}{(\alpha+\beta)^2(\alpha+\beta+1)}$ \\
\bottomrule
\end{tabular}
\end{center}

\begin{recoverybox}
Before the exercises, the one-sentence summary of each story. \emph{Uniform}:
equally likely across an interval -- pure ignorance, variance $(b-a)^2/12$.
\emph{Exponential}: memoryless waiting time, constant hazard, mean $=$ SD $=
1/\lambda$. \emph{Gamma}: the wait for the $\alpha$-th event, a sum of
exponentials, mean $\alpha/\lambda$. \emph{Beta}: a flexible distribution over a
probability on $[0,1]$, the conjugate prior whose update is ``add heads to
$\alpha$, tails to $\beta$.'' Four shapes, one connected family.
\end{recoverybox}

\section{Exercises}
\label{ch11:exercises}

\begin{exercise}[Uniform basics]
\label{ch11:exr-unif}
Let $X\sim\Unif(2,8)$. (a) Write down the density and CDF. (b) Compute
$\P(3 \le X \le 5)$. (c) Find $\E[X]$ and $\Var(X)$. (d) Find the value $m$ with
$\P(X \le m) = 0.9$.
\end{exercise}

\begin{exercise}[The quantization-noise rule]
\label{ch11:exr-quant}
An analog-to-digital converter rounds voltages to the nearest step of size
$\Delta$. Model the rounding error as $\Unif(-\Delta/2,\Delta/2)$. (a) Show the
error has mean $0$ and variance $\Delta^2/12$. (b) If the signal itself has power
(variance) $\sigma^2$, the signal-to-noise ratio is $\sigma^2/(\Delta^2/12)$.
Explain in words why halving the step size $\Delta$ improves the SNR by a factor
of $4$ (about $6$ decibels).
\end{exercise}

\begin{exercise}[Exponential computations]
\label{ch11:exr-expo}
Calls arrive at a help desk at rate $\lambda = 4$ per hour with exponential
inter-arrival times. (a) What is the probability the next call takes more than
$30$ minutes to arrive? (b) What is the expected wait, and the median wait?
(c) Given that no call has come in the first $20$ minutes, what is the expected
\emph{additional} wait? Justify your answer with the memoryless property.
\end{exercise}

\begin{exercise}[Memorylessness, proved both ways]
\label{ch11:exr-memoryless}
(a) Using only $\P(T>t)=e^{-\lambda t}$, verify directly that
$\P(T > s+t \given T > s) = \P(T>t)$. (b) Conversely, suppose a continuous
positive variable $T$ satisfies $g(s+t)=g(s)g(t)$ for $g(t)=\P(T>t)$, all
$s,t\ge0$, with $g$ decreasing and $g(0)=1$. By considering $g(t)$ for rational
$t$ first, argue that $g(t)=e^{-\lambda t}$ for some $\lambda>0$, so the
exponential is the \emph{only} memoryless continuous distribution.
\end{exercise}

\begin{exercise}[Minimum of exponentials and who wins the race]
\label{ch11:exr-min}
Let $T_1\sim\Expo(\lambda_1)$ and $T_2\sim\Expo(\lambda_2)$ be independent.
(a) Show $\min(T_1,T_2)\sim\Expo(\lambda_1+\lambda_2)$. (b) Show that the
probability that component $1$ fails first is
$\P(T_1 < T_2) = \lambda_1/(\lambda_1+\lambda_2)$. \emph{Hint: condition on
$T_1$ and integrate, or set up the double integral over $\{t_1<t_2\}$.}
(c) Interpret (b): why is the faster component more likely to win?
\end{exercise}

\begin{exercise}[Gamma as a sum]
\label{ch11:exr-gamma}
Machine parts arrive at rate $\lambda=3$ per hour, exponential gaps. (a) Name
the distribution of the wait until the $4$th part. (b) Give its mean, variance,
and standard deviation. (c) Compare the ratio SD/mean for this wait against the
ratio for a single arrival, and explain why summing waits reduces relative
variability.
\end{exercise}

\begin{exercise}[Verifying a gamma normalizing constant]
\label{ch11:exr-gammaconst}
Using the substitution $u=\lambda x$ and the definition
$\Gamma(\alpha)=\int_0^\infty u^{\alpha-1}e^{-u}\dd u$, show that
$\int_0^\infty \frac{\lambda^\alpha}{\Gamma(\alpha)} x^{\alpha-1}e^{-\lambda x}\dd x = 1$.
Then use $\Gamma(\alpha+1)=\alpha\Gamma(\alpha)$ to derive $\E[X]=\alpha/\lambda$
directly from the integral $\int_0^\infty x\,f(x)\,\dd x$.
\end{exercise}

\begin{exercise}[Reading and updating a beta]
\label{ch11:exr-beta}
You are testing a new web button. Your prior on its click-through rate $p$ is
$\Bet(2,8)$. (a) What is your prior mean and standard deviation for $p$? (b) You
show the button to $40$ visitors; $12$ click. Write down the posterior
distribution of $p$. (c) Give the posterior mean. Is it closer to the prior mean
or to the raw data proportion $12/40$? Explain why.
\end{exercise}

\begin{exercise}[The uniform is a beta]
\label{ch11:exr-unifbeta}
(a) Show that $\Bet(1,1)$ has density equal to $1$ on $[0,1]$, i.e.\ it is
$\Unif(0,1)$. (b) Using the general beta mean and variance formulas, recover the
$\Unif(0,1)$ mean $1/2$ and variance $1/12$. (c) Conceptually: explain why
starting Bayesian updating from a $\Bet(1,1)$ prior corresponds to ``assuming
nothing'' about the unknown probability.
\end{exercise}

\begin{exercise}[The length-biased waiting paradox]
\label{ch11:exr-paradox}
Buses have i.i.d.\ exponential gaps with mean $10$ minutes. (a) Using
memorylessness, explain why a passenger arriving at a random time waits, on
average, $10$ minutes for the next bus -- not $5$. (b) Argue that the time since
the previous bus also averages $10$ minutes, so the observed gap averages $20$.
(c) In one or two sentences, explain the resolution: why does sampling by a
random \emph{instant} over-represent long gaps?
\end{exercise}

\begin{exercise}[Proof: scaling an exponential]
\label{ch11:exr-scale}
Let $T\sim\Expo(\lambda)$ and let $c>0$. Using CDFs, prove that
$cT\sim\Expo(\lambda/c)$. Interpret: rescaling time (e.g.\ measuring in seconds
instead of minutes) just rescales the rate, as it must.
\end{exercise}

\begin{exercise}[Inverse-transform simulation, by hand]
\label{ch11:exr-invtransform}
You have a uniform draw $U=0.2$. (a) Using $T=-\tfrac{1}{\lambda}\ln(1-U)$ with
$\lambda=2$, compute the corresponding exponential sample. (b) Explain why
$-\tfrac{1}{\lambda}\ln(1-U)$ and $-\tfrac{1}{\lambda}\ln U$ give samples from
the \emph{same} distribution. (c) Why does the inverse-transform method turn a
flat uniform into a right-skewed exponential -- what is the logarithm doing
geometrically to the unit interval?
\end{exercise}
